\subsection{Программная реализация ключевых компонентов}

В данном разделе представлены основные программные компоненты системы анализа обитаемости экзопланет. Каждый компонент реализует определённую функциональность и вносит свой вклад в общую работу системы.

\subsubsection{Валидация входных параметров}

Одним из ключевых аспектов надёжной работы системы является корректная валидация входных параметров экзопланеты. Листинг \ref{lst:validation} демонстрирует реализацию функции проверки параметров.

\begin{lstlisting}[language=JavaScript, caption={Функция валидации параметров экзопланеты}, label={lst:validation}]
function validatePlanetParameters(radius, temperature, density) {
    const validation = {
        isValid: true,
        errors: []
    };
    
    // Проверка радиуса (0.1 - 3.0 радиусов Земли)
    if (radius < 0.1 || radius > 3.0) {
        validation.isValid = false;
        validation.errors.push('Радиус должен быть в диапазоне 0.1-3.0 радиусов Земли');
    }
    
    // Проверка температуры (50K - 1000K)
    if (temperature < 50 || temperature > 1000) {
        validation.isValid = false;
        validation.errors.push('Температура должна быть в диапазоне 50-1000K');
    }
    
    // Проверка плотности (0.01 - 3.0 относительно Земли)
    if (density < 0.01 || density > 3.0) {
        validation.isValid = false;
        validation.errors.push('Плотность должна быть в диапазоне 0.01-3.0 плотности Земли');
    }
    
    return validation;
}
\end{lstlisting}

Функция \texttt{validatePlanetParameters} принимает три основных параметра экзопланеты:
\begin{itemize}
    \item radius — радиус планеты относительно земного (0.1--3.0 R⊕)
    \item temperature — температура поверхности в Кельвинах (50--1000 K)
    \item density — плотность относительно земной (0.01--3.0 ρ⊕)
\end{itemize}

Каждый параметр проверяется на соответствие допустимому диапазону значений, определённому на основе актуальных астрофизических данных и ограничений современных методов наблюдения экзопланет.

\subsubsection{Расчёт индекса подобия Земле}

Индекс подобия Земле (Earth Similarity Index, ESI) является ключевым показателем при оценке потенциальной обитаемости экзопланеты. Листинг \ref{lst:esi} представляет реализацию алгоритма расчёта ESI.

\begin{lstlisting}[language=JavaScript, caption={Функция расчёта индекса подобия Земле}, label={lst:esi}]
function calculateESI(radius, temperature, density) {
    // Параметры Земли для сравнения
    const EARTH_TEMP = 288; // K
    const EARTH_RADIUS = 1.0; // относительные единицы
    const EARTH_DENSITY = 1.0; // относительные единицы
    
    // Весовые коэффициенты для параметров
    const TEMP_WEIGHT = 0.5;
    const RADIUS_WEIGHT = 0.3;
    const DENSITY_WEIGHT = 0.2;
    
    // Расчет компонентов ESI
    const tempESI = Math.pow((1 - Math.abs(temperature - EARTH_TEMP) / 
        (temperature + EARTH_TEMP)), TEMP_WEIGHT);
    const radiusESI = Math.pow((1 - Math.abs(radius - EARTH_RADIUS) / 
        (radius + EARTH_RADIUS)), RADIUS_WEIGHT);
    const densityESI = Math.pow((1 - Math.abs(density - EARTH_DENSITY) / 
        (density + EARTH_DENSITY)), DENSITY_WEIGHT);
    
    // Итоговый ESI
    const ESI = tempESI * radiusESI * densityESI;
    
    return ESI;
}
\end{lstlisting}

Алгоритм расчёта ESI учитывает три ключевых параметра с соответствующими весовыми коэффициентами:
\begin{itemize}
    \item Температура (вес 0.5) — наиболее значимый параметр для оценки обитаемости
    \item Радиус (вес 0.3) — определяет силу гравитации и способность удерживать атмосферу
    \item Плотность (вес 0.2) — влияет на состав и структуру планеты
\end{itemize}

\subsubsection{Визуализация 3D-модели экзопланеты}

Для наглядного представления экзопланеты реализована интерактивная 3D-модель с использованием библиотеки Three.js. Листинг \ref{lst:3dmodel} демонстрирует код инициализации модели.

\begin{lstlisting}[language=JavaScript, caption={Инициализация 3D-модели экзопланеты}, label={lst:3dmodel}]
function initPlanet3DModel(radius, temperature) {
    const scene = new THREE.Scene();
    const camera = new THREE.PerspectiveCamera(75, 
        window.innerWidth / window.innerHeight, 0.1, 1000);
    const renderer = new THREE.WebGLRenderer({ antialias: true });
    
    // Создание геометрии планеты
    const geometry = new THREE.SphereGeometry(radius, 32, 32);
    
    // Материал планеты в зависимости от температуры
    const material = new THREE.MeshPhongMaterial({
        map: generatePlanetTexture(temperature),
        bumpMap: generateBumpMap(),
        bumpScale: 0.05,
        specular: new THREE.Color('grey'),
        shininess: 10
    });
    
    const planet = new THREE.Mesh(geometry, material);
    scene.add(planet);
    
    // Добавление освещения
    const light = new THREE.PointLight(0xffffff, 1, 100);
    light.position.set(10, 10, 10);
    scene.add(light);
    
    // Настройка орбитального управления
    const controls = new THREE.OrbitControls(camera, renderer.domElement);
    controls.enableDamping = true;
    controls.dampingFactor = 0.05;
    
    return { scene, camera, renderer, planet, controls };
}
\end{lstlisting}

3D-визуализация включает следующие ключевые компоненты:
\begin{itemize}
    \item Создание сферической геометрии с учётом радиуса планеты
    \item Генерация текстуры поверхности на основе температуры
    \item Настройка освещения для реалистичного отображения
    \item Реализация интерактивного управления камерой
\end{itemize}

\subsubsection{Визуализация сравнительного анализа}

Для сравнения параметров экзопланеты с земными реализована система построения графиков с использованием библиотеки Plotly. Листинг \ref{lst:charts} показывает код генерации сравнительных диаграмм.

\begin{lstlisting}[language=JavaScript, caption={Генерация графиков сравнения с Землёй}, label={lst:charts}]
function generateComparisonCharts(planetData) {
    const trace1 = {
        type: 'scatterpolar',
        r: [
            planetData.radius,
            planetData.temperature / 288, // нормализация к температуре Земли
            planetData.density,
            planetData.esi
        ],
        theta: ['Радиус', 'Температура', 'Плотность', 'ESI'],
        fill: 'toself',
        name: 'Экзопланета'
    };
    
    const trace2 = {
        type: 'scatterpolar',
        r: [1, 1, 1, 1], // параметры Земли
        theta: ['Радиус', 'Температура', 'Плотность', 'ESI'],
        fill: 'toself',
        name: 'Земля'
    };
    
    const layout = {
        polar: {
            radialaxis: {
                visible: true,
                range: [0, 2]
            }
        },
        showlegend: true
    };
    
    Plotly.newPlot('comparison-chart', [trace1, trace2], layout);
}
\end{lstlisting}

Система визуализации обеспечивает:
\begin{itemize}
    \item Построение полярных диаграмм для сравнения параметров
    \item Нормализацию значений относительно земных показателей
    \item Интерактивное взаимодействие с графиками
    \item Наглядное представление различий между планетами
\end{itemize}

\subsubsection{Анализ вероятности наличия жидкой воды}

Оценка вероятности наличия жидкой воды является критически важным компонентом анализа обитаемости. Листинг \ref{lst:water} демонстрирует реализацию алгоритма расчёта.

\begin{lstlisting}[language=JavaScript, caption={Расчёт вероятности наличия жидкой воды}, label={lst:water}]
function calculateWaterProbability(temperature, radius, density) {
    // Константы для расчета
    const WATER_FREEZING_POINT = 273.15; // K
    const WATER_BOILING_POINT = 373.15; // K
    const GRAVITY_FACTOR = Math.sqrt(density * Math.pow(radius, 3));
    
    // Корректировка точек фазового перехода с учетом гравитации
    const adjustedFreezingPoint = WATER_FREEZING_POINT * 
        Math.pow(GRAVITY_FACTOR, 0.1);
    const adjustedBoilingPoint = WATER_BOILING_POINT * 
        Math.pow(GRAVITY_FACTOR, 0.15);
    
    // Расчет вероятности
    let probability = 0;
    
    if (temperature >= adjustedFreezingPoint && 
        temperature <= adjustedBoilingPoint) {
        // Максимальная вероятность в середине диапазона
        const optimalTemp = (adjustedFreezingPoint + adjustedBoilingPoint) / 2;
        const deviation = Math.abs(temperature - optimalTemp);
        const maxDeviation = (adjustedBoilingPoint - adjustedFreezingPoint) / 2;
        
        probability = 1 - (deviation / maxDeviation);
    }
    
    return probability;
}
\end{lstlisting}

Алгоритм учитывает следующие факторы:
\begin{itemize}
    \item Температуру поверхности планеты
    \item Гравитационное влияние на точки фазового перехода воды
    \item Оптимальный диапазон температур для жидкой воды
    \item Корректировку на основе массы и плотности планеты
\end{itemize}

Представленные программные компоненты формируют основу системы анализа обитаемости экзопланет, обеспечивая точные расчёты, наглядную визуализацию и удобный пользовательский интерфейс. Каждый компонент разработан с учётом современных научных данных и требований к точности астрофизических расчётов. 